\section{Conclusion}

\ifshort
For exfiltration the per-episode judgement can be replaced by a construction:
plant a record no legitimate task transmits, and read its bytes at the egress
boundary. The result worth carrying forward required admitting a design error---
the obvious form of that label had a \llamaControlFPRate{} false-positive rate on
traffic with no adversary in it.
\else

Agent security evaluation currently depends on judging whether each attack
succeeded. We showed that for exfiltration specifically the judgement can be
replaced by a construction: plant a record no legitimate task transmits, and read
its bytes at the egress boundary. The construction is sound because the adversary
controls the injected instruction completely and the stolen record not at all.

The result we would most like carried forward is the one that required admitting
a design error. The obvious form of this label---the planted record in any
outbound argument---had a \llamaControlFPRate{} false-positive rate on
adversary-free traffic, because transmitting member records is a plausible
reading of the agent's job. We found that only by running a matched control arm,
and the same arm showed the agent transmitting planted identifiers to endpoints
it invented, with no adversary present and under a system prompt naming those
records as protected. Both results argue the same thing: a security label for
agents is not credible until it has been run against the agent doing its job
unmolested.
\fi
